\subsection{Identificação paramétrica e engenharia reversa}

A aplicação direta dos coeficientes apresentados por Altıntaş e
Ertan~\cite{altintas2024} não permitiu reproduzir simultaneamente o nível,
a inclinação e a curvatura das curvas publicadas. Além disso, o artigo não
explicita de forma inequívoca o número de células do empilhamento, a área
ativa, a expressão utilizada para a eficiência, a temperatura do ensaio de
pressão e algumas convenções dimensionais das equações. Dessa forma, foi
realizada uma etapa de engenharia reversa seguida de identificação
paramétrica, preservando-se a estrutura eletroquímica do modelo e tratando
os valores não fornecidos como parâmetros inferidos.

Inicialmente, as quatro curvas de tensão da Fig.~3 foram digitalizadas e os
eixos foram recalibrados diretamente a partir da imagem original. Foram
consideradas as curvas correspondentes a $T=298{,}15~\mathrm{K}$,
$T=373{,}15~\mathrm{K}$, $P_{ar}=1~\mathrm{atm}$ e
$P_{ar}=5~\mathrm{atm}$. As curvas de potência e eficiência não foram
incluídas na função objetivo principal, uma vez que ambas são derivadas da
tensão e sua inclusão produziria dupla contagem da mesma informação.

O vetor de parâmetros identificado foi definido por

\begin{equation}
\boldsymbol{\theta}=
\left[
\xi_1,\xi_2,\xi_3,\xi_4,
R_{\mathrm{mem,ref}},n_R,
a_{\mathrm{ref}},a_T,b
\right]^{\mathrm{T}}.
\end{equation}

Os quatro primeiros valores iniciais foram obtidos diretamente da equação
de ativação publicada. A resistência de referência foi inicialmente
estimada a partir da inclinação da região aproximadamente ôhmica, enquanto
os parâmetros de concentração foram inicializados com base na expressão
apresentada no artigo e na respectiva conversão de unidades. Para reduzir a
dependência do valor inicial, foram utilizados 16 inícios determinísticos,
gerados dentro de limites fisicamente admissíveis com semente fixa igual a
20260720.

A identificação foi formulada como um problema de mínimos quadrados não
lineares limitado:

\begin{equation}
J(\boldsymbol{\theta})=
\frac{1}{N}
\sum_{c=1}^{4}
\sum_{k=1}^{n_c}
\left[
\frac{
V_{\mathrm{mod}}(j_k,T_c,P_c;\boldsymbol{\theta})-
V_{\mathrm{fig},c}(j_k)
}{\sigma_V}
\right]^2,
\end{equation}

em que $V_{\mathrm{mod}}$ representa a tensão calculada pelo modelo,
$V_{\mathrm{fig}}$ corresponde ao valor digitalizado e
$\sigma_V=0{,}003~\mathrm{V}$ representa a incerteza adotada para a
extração gráfica. A minimização foi realizada pelo método Trust Region
Reflective, implementado na função \texttt{least\_squares} da biblioteca
SciPy, com limites inferiores e superiores para todos os parâmetros. O
conjunto final foi selecionado pelo menor erro quadrático médio conjunto das
quatro curvas de tensão.

Após a identificação eletroquímica, os parâmetros escalares do empilhamento
foram inferidos separadamente. O produto $N A$ foi estimado a partir de

\begin{equation}
P_{\mathrm{stack}}=N A j V_{\mathrm{cell}}.
\end{equation}

A área ativa de $232~\mathrm{cm^2}$ foi associada ao termo de capacitância
$0{,}035\times232~\mathrm{F}$ apresentado na tabela de parâmetros. O ajuste
contínuo das curvas de potência resultou em um número equivalente de
aproximadamente $41{,}70$ células. Como o número de células deve ser inteiro,
foi adotado $N=41$, valor que preserva a ordem de grandeza de potência e a
consistência com a curva de maior temperatura. A diferença residual entre as
curvas de potência foi mantida e documentada, sem a introdução de fatores de
correção ocultos.

A expressão de eficiência não é fornecida no artigo. A relação observada
entre tensão e eficiência foi representada por

\begin{equation}
\eta=100\,U_f\frac{V_{\mathrm{cell}}}{E_0},
\end{equation}

sendo $U_f$ determinado por mínimos quadrados escalares a partir das duas
curvas de eficiência. Por fim, a identificabilidade local foi avaliada por
meio do Jacobiano escalado do problema, incluindo o número de condição e a
matriz de correlação entre os parâmetros. Esse procedimento permite
distinguir a capacidade de reprodução das curvas da interpretação física
individual dos coeficientes efetivos.
